\section*{Limitations}
\label{sec:limitations}

Our evaluation is restricted to multiple-choice benchmarks; the flip metric may not transfer to open-ended settings, where sycophancy is more likely to show up as semantic drift than a discrete answer change. HLE sample sizes are small ($24$--$72$ vs.\ $400$--$420$ $T_0$-correct questions per condition), so only one model's slope reaches significance individually, though the direction is uniform across all measurable conditions. The confidence estimator introduces its own error into the signal being measured; we see no qualitatively different miscalibration there (\autoref{app:calibration}). The planted wrong answer $a^\dagger$ is drawn uniformly rather than targeted at the model's most plausible error. Using the existing baseline samples, we find flip likelihood does increase with $a^\dagger$'s baseline plausibility once model is controlled for ($\hat\beta=+0.85$, $p<10^{-6}$; \autoref{app:plausibility}), so reported flip rates likely understate what a genuinely targeted adversary could achieve. Our pressure sequence uses one fixed category order; the single-turn experiment rules out the simplest version of this concern, but a category-by-position interaction is not ruled out. Our six prompts are pure rhetorical pressure, so the flip metric cannot distinguish capitulation from legitimate belief updating.